% from ega3-7-fr.tex
\oldpage[III]{197}
En effet, partant de la suite exacte
$0\to Z_1\to P_1\to B_0\to0$, on en tire la suite exacte
\[
  0\longrightarrow\operatorname{Tor}_1^A(B_0,k)
  \longrightarrow Z_1\otimes k
  \longrightarrow P_1\otimes k
  \xrightarrow{u}B_0\otimes k
  \longrightarrow0.
\]
De la suite exacte $0\to B_0\to Z_0\to Z'_0\to0$, on tire, puisque $Z_0=P_0$
est plat,
$\operatorname{Tor}_1^A(B_0,k)=\operatorname{Tor}_2^A(Z'_0,k)$; par définition,
on a $Z_1=T_1(A)$; enfin, on a $T_1(k)=\operatorname{Ker}(d\otimes1)$, et
$d\otimes1$ se factorise en
\[
  P_1\otimes k
  \xrightarrow{u}B_0\otimes k
  \xrightarrow{v}Z_0\otimes k
  =P_0\otimes k;
\]
on a $T_1(k)=u^{-1}(R)$, où $R=\operatorname{Ker}v$, et comme $u$ est
surjectif, $R=u(T_1(k))$; enfin,
$R=\operatorname{Tor}_1^A(Z'_0,k)$ par définition de $v$, ce qui achève
d'établir la suite exacte (7.6.11.7).
\end{proof}

Pour déduire alors c) de b), on considère une suite exacte
\[
  L_1\xrightarrow{d}L_0\longrightarrow N\longrightarrow0,
\]
où $L_0$ et $L_1$ sont des modules libres de type fini; considérons le foncteur
$T$ associé au complexe formé de $L_1$ et $L_0$; comme $L_0$ et $L_1$ sont
libres, ils s'identifient à leurs biduals; donc si
\[
  M=\operatorname{Coker}({}^t d),
  \qquad
  T_1(A)=\operatorname{Ker}(d)=\check M;
\]
par ailleurs, $M\otimes_A k=\operatorname{Coker}({}^t d\otimes1_k)$ a même rang
sur $k$ que $\operatorname{Ker}(d\otimes1_k)$. L'hypothèse b) entraîne par suite
que
\[
  \operatorname{rang}_k\bigl(T_1(A)\otimes_A k\bigr)
  \leq
  \operatorname{rang}_k\bigl(T_1(k)\bigr);
\]
comme $Z'_0=N$, l'inégalité (7.6.11.4) résulte donc de la suite exacte
(7.6.11.7).

Enfin, pour prouver que c) entraîne a), appliquons (7.6.11.6) en remplaçant
$P_0$ et $P_1$ par $P_p$ et $P_{p-1}$; l'hypothèse c) appliquée au module
$Z'_p$ donne $\operatorname{rang}_kR\geq\operatorname{rang}_kS$, où
\[
  R=\operatorname{Ker}(P_p\otimes k\longrightarrow P_{p-1}\otimes k)
  \quad\text{et}\quad
  S=Z_p\otimes k.
\]
Or, si l'on factorise $d_{p+1}:P_{p+1}\to P_p$ en
\[
  P_{p+1}\xrightarrow{v}Z_p\xrightarrow{j}P_p,
\]
on a
\[
  \operatorname{Im}(d_{p+1}\otimes1)
  =
  (j\otimes1)\bigl(\operatorname{Im}(v\otimes1)\bigr).
\]
Comme
\[
  T_p(A)\otimes_A k
  =(Z_p/B_p)\otimes_A k
  =(Z_p\otimes k)/\operatorname{Im}(v\otimes1),
\]
et
\[
  T_p(k)=R/\operatorname{Im}(d_{p+1}\otimes1),
\]
on en conclut bien l'inégalité (7.6.11.2).

Cela étant, supposons que l'anneau local $A$ soit régulier de dimension $n$;
on sait alors [17] que le $A$-module
$\operatorname{Tor}_i^A(k,k)$ est isomorphe à la puissance extérieure
\[
  \bigwedge\nolimits^i(\mathfrak m/\mathfrak m^2);
\]
on voit donc que la condition (7.6.11.4) n'est pas vérifiée pour $N=k$, dès que
$n\geq4$.

Par contre, si l'anneau local intègre $A$ est tel que tout $A$-module réflexif
de type fini soit libre (ce qui est le cas lorsque $A$ est un anneau régulier
de dimension 2), la condition (7.6.11.3) est vérifiée : en effet, on sait que le
dual $\check M$ d'un $A$-module $M$ de type fini est réflexif, donc libre, et
par suite
\[
  \operatorname{rang}_k(\check M\otimes_A k)
  =
  \operatorname{rang}_K(\check M)
  =
  \operatorname{rang}_K(M),
\]
($K$ corps des fractions de $A$); par ailleurs, on sait que toute base sur $k$
de $M\otimes_A k$ est formée d'images d'un système de générateurs de $M$
(Bourbaki, \emph{Alg. comm.}, chap. II, \S 3, n\textsuperscript{o} 2, cor. 2 de
la prop. 4), donc
$\operatorname{rang}_K(M)\leq\operatorname{rang}_k(M\otimes_A k)$, ce qui
prouve notre assertion.
\end{remark}

\subsection{Application aux morphismes propres : I. La propriété d'échange}
\label{subsection:III.7.7.fr}

Les trois numéros qui suivent sont, pour l'essentiel, des traductions, dans le
langage des morphismes de préschémas, des résultats des numéros précédents.

\begin{env}[7.7.1]
Soit $f:X\to Y$ un morphisme quasi-compact et séparé de préschémas, et soit
$\mathscr P_\bullet$ un complexe de $\mathscr O_X$-Modules quasi-cohérents, dont
l'opérateur de dérivation est de degré $-1$; supposons en outre que les
$\mathscr O_X$-Modules $\mathscr P_i$ sont $Y$-plats ($0_{\mathrm I}$, 6.7.1).
